\section{Introduction}

\subsection{Problem Statement and Motivation}
% [REVISION: Adjusted to Supervisor's 'Funnel' structure: AI -> Epilepsy -> Context]

% 1. Broad Context: AI in Medicine
Artificial intelligence (AI) has made substantial advances in medical
signal analysis over the last decade, enabling new possibilities for
continuous monitoring and early diagnostics. Deep learning, in
particular, has shifted the paradigm from manual feature engineering to
automated representation learning, showing high efficacy in
interpreting complex physiological time-series data. These technological
strides offer a pathway to support clinical decision-making and
automate patient monitoring in real-time environments.

% 2. Specific Application: Epilepsy
A critical application domain for these technologies is epilepsy, a
neurological disorder affecting roughly 7.6 per 1,000 people globally
\parencite{beghiEpidemiologyEpilepsy2019}. With approximately 30\%
of patients being drug-resistant, there is an urgent demand for timely,
low-burden monitoring systems that can facilitate early
intervention \parencite{kwanEarlyIdentificationRefractory2000}. While
Electroencephalography (EEG) remains the gold standard for seizure
detection in controlled clinical settings, its hardware is too
obtrusive for long-term pervasive use in daily life.

% 3. The Gap/Challenge: Clinic vs. Ambulatory & EEG vs. Non-EEG
This creates a divergence between clinical feasibility and ambulatory
needs. Wearable devices measuring non-EEG autonomic biosignals—such as
heart rate (ECG/PPG), electrodermal activity (EDA), and
accelerometry (ACC)—have emerged as a patient-friendly alternative.
However, applying deep learning to these signals introduces distinct
challenges regarding reliability. The central problem is
differentiating between the high-precision requirements of clinical
diagnostics and the high-sensitivity, low-false-alarm requirements
necessary for ambulatory safety.

\subsection{Research Aim and Questions}
The aim of this literature review is to analyze machine and deep
learning models for seizure detection and prediction using non-EEG
wearable biosignals. Specifically, it seeks to identify which modeling
strategies perform optimally in ambulatory scenarios compared to
controlled clinical environments.

\subsubsection{Primary Research Question}
In non-EEG wearable seizure detection and prediction
(ECG/HRV, PPG, EDA, ACC), which modeling and evaluation choices
reported in 2015--2025 are associated with clinically viable
performance (low false-alarm rates and high sensitivity) in ambulatory
settings?

\subsubsection{Sub-questions}
\begin{enumerate}
  \item What ranges of sensitivity, FAR/h, and latency are reported under
    retrospective, pseudo-prospective, and prospective designs?
  \item How do traditional feature-based models compare to deep temporal
    architectures (1D-CNN, LSTM, TCN, Transformers) regarding
    computational load and latency?
  \item What evidence exists regarding robustness to domain shift
    (patient/device variability) and how do studies address the
    "Black Box" interpretability vs. performance trade-off?
\end{enumerate}

\subsection{Methodological Approach and Structure}
This review follows structured literature review guidance according to
\cite{websterAnalyzingPastDetermine2002} and adheres to PRISMA principles
for transparent reporting.

The search strategy targets the period 2015--2025 using a multi-stage
approach:
\begin{enumerate}
  \item \textbf{Database Search:} Structured queries in IEEE Xplore,
    PubMed, and Scopus to identify peer-reviewed studies on wearable
    seizure detection.
  \item \textbf{Screening:} A two-phase screening (title/abstract and
    full-text) to exclude EEG-only studies and purely simulation-based
    papers.
  \item \textbf{Analysis:} Extracted literature will be synthesized
    using a concept matrix to categorize studies by model
    architecture, study phase (clinical vs. ambulatory), and
    evaluation metrics.
\end{enumerate}

\section{Theoretical Background and Clinical Concepts}

% [REVISION: Added this section to address Supervisor's point on "Healthcare specifics"]
\subsection{Clinical Deployment Requirements}
Deploying AI in healthcare requires balancing competing performance
metrics. In epilepsy monitoring, \textbf{sensitivity} (detecting all
seizures to ensure safety) must be weighed against \textbf{precision}
(minimizing false alarms). In a clinical ward, high sensitivity is
prioritized as medical staff can verify alarms. In ambulatory (home)
settings, however, a high False Alarm Rate (FAR) causes "alarm
fatigue," rendering the device useless for the patient.

Furthermore, the "Black Box" nature of deep learning poses safety
challenges. Unlike traditional decision trees, deep neural networks
often lack transparency. For clinical acceptance, systems generally
require explainability to support, rather than replace, physician
judgment.

\subsection{Epileptic Seizures and Physiology}
Epileptic seizures are classified by onset as focal or generalized
\parencite{fisherInstructionManualILAE2017}. While EEG captures the
electrical brain onset, the autonomic nervous system provides surrogate
markers. Sympathetic activation manifests as ictal tachycardia,
reduced Heart Rate Variability (HRV), and increased Electrodermal
Activity (EDA), which serve as the inputs for the models reviewed in
this work.

\section{Modeling Landscape and Design Dimensions}
(This section remains as previously drafted, focusing on CNN/LSTM/TCN,
Fusion, and Personalization...)

\newpage
\section*{Provisional Table of Contents}
\textit{Target Length: Approx. 15 Pages of Text}

\begin{enumerate}
  \item **Introduction** (approx. 2.5 pages)
    \begin{itemize}
      \item 1.1 Problem Statement and Motivation (AI advances $\rightarrow$ Epilepsy context)
      \item 1.2 Aim and Research Questions
      \item 1.3 Methodological Approach (Webster \& Watson / PRISMA)
    \end{itemize}
    
  \item **Theoretical Background and Clinical Concepts** (approx. 2.5 pages)
    \begin{itemize}
      \item 2.1 Epilepsy and Autonomic Biosignals (EEG vs. Non-EEG)
      \item 2.2 Clinical Performance Requirements (Sensitivity vs. Precision, FAR/h)
      \item 2.3 Safety, Interpretability, and the "Black Box" Problem
      \item 2.4 Study Phases: Retrospective vs. Prospective
    \end{itemize}
    
  \item **Modeling Landscape** (approx. 6 pages)
    \begin{itemize}
      \item 3.1 Feature Engineering vs. Deep Learning (CNN/LSTM/TCN/Transformers)
      \item 3.2 Sensor Modalities and Fusion Strategies
      \item 3.3 Personalization and Domain Shift
    \end{itemize}
    
  \item **Results and Discussion** (approx. 3 pages)
    \begin{itemize}
      \item 4.1 Comparative Performance (Clinical vs. Ambulatory)
      \item 4.2 Trade-offs: Complexity vs. Latency
    \end{itemize}
    
  \item **Conclusion and Future Work** (approx. 1 page)
\end{enumerate}